\emph{(i).} If $\langle\cdot,\cdot\rangle$ is positive definite, $Q^{*} = Q$ and $Q^2 = 0$, then $\langle Qx, Qx\rangle = \langle x, Q^2x\rangle = 0$ for every $x$, so $Qx = 0$. On a floor of the tower (Fact \claimref{field-theory.imported-tower-limit}) the form $\langle x^{*}y\rangle_0$, with $*$ the reversion followed by a sign for each negative-square generator, makes the basis blades orthonormal: the check confirms $\langle b^{*}c\rangle_0 = \delta_{bc}$ over all blades of floors with one and two planes, and the computation for a general blade is the same product of signs, each factor $e^{*}e = e^2 = 1$ or $f^{*}f = -f^2 = 1$. On the palatial polynomial representation of Fact \claimref{field-theory.imported-statistics-no-go-palatial-algebra} the form $\langle z^n, z^m\rangle = n!\,\delta_{nm}$ is diagonal with positive entries. So a nonzero nilpotent charge is not self-adjoint for either form, and the self-adjoint charge of the Kugo-Ojima framework, whose physical states are $\ker Q/\operatorname{im}Q$ with $Q^{\dagger} = Q$, does not exist on them. Nilpotent charges that are not self-adjoint exist on positive spaces (their cohomology is computed with $Q$ and $Q^{*}$), so nothing here excludes covariant gauge fixing in every form; the general statement that a local covariant vector potential obeying Maxwell's equations needs an indefinite metric is Strocchi's theorem (Fact \claimref{synthesis.imported-brst-kugo-ojima-strocchi-and-gribov}(iii)); the check exhibits this for the charge of (ii), which is nilpotent and nonzero and fails $PQ = Q^{\mathsf T}P$ for the positive Gram matrix $P$.

\emph{(ii).} With $a_\mu = \partial/\partial z_\mu$ and positive adjoint $a_\mu^{*} = z_\mu$, the operator $(-1)^{n_0}$ is an involution that anticommutes with $z_0$ and commutes with $z_i$, so the Krein adjoint $a_\mu^{\dagger} = \eta a_\mu^{*}\eta$ is $-z_0$ for $\mu = 0$ and $z_i$ for $\mu = i$, that is $-g_{\mu\mu}z_\mu$. Hence $[a_\mu, a_\nu^{\dagger}] = -g_{\mu\nu}$, and for $a(u) = \sum_\mu g_{\mu\mu}u^\mu a_\mu$ one gets $[a(u), a(w)^{\dagger}] = -u\cdot w$, the Minkowski metric of the polarisations with a sign: the standard Gupta-Bleuler form, the same in any formalism. A boost rotor $R$ preserves $u\cdot w$, so this one-particle form is boost-invariant; the check confirms it for $R = \tfrac54 + \tfrac34\gamma_0\gamma_1$ ($R\tilde R = 1$) and shows the Euclidean form $\sum_\mu u^\mu w^\mu$ changes. On the ghost floor, $S = 1 - (b^{*} - d^{*})(b - d)$ is an element of the floor with $S^2 = 1$, $S^{*} = S$ and $SbS = d$ (checked), so $b^{\dagger} = Sb^{*}S = d^{*}$ and $(d^{*})^{\dagger} = b$. Bosons and ghosts commute. For $L = \sum_\mu g_{\mu\mu}k^\mu a_\mu$, $L^{\dagger} = -\sum_\mu k^\mu z_\mu$ and $[L, L^{\dagger}] = -Q(k) = 0$. Then
\[
Q^2 = L^{\dagger}L\,(bd^{*} + d^{*}b) + (L^{\dagger})^2b^2 + d^{*2}L^2 + [L, L^{\dagger}]\,d^{*}b = 0,
\]
since $b$ and $d^{*}$ anticommute, $b^2 = d^{*2} = 0$ and $k$ is null; and $Q^{\dagger} = b^{\dagger}L + L^{\dagger}d = d^{*}L + L^{\dagger}b = Q$, the bosonic and ghost factors commuting. The check confirms $Q^2 = 0$, $KQ = Q^{\mathsf T}K$ and $K = K^{\mathsf T}$ for the Krein Gram matrix $K = P\eta$ on the space with $T \le 3$ (which $Q$ preserves; in the truncation the factor that lowers $T$ is applied first), for $k = \gamma_0 + \gamma_3$ and $k = 5\gamma_0 + 3\gamma_1 + 4\gamma_2$.

\emph{(iii).} $Q$ commutes with $T = \deg + N_b + N_d$ and raises the ghost number $N_d - N_b$ by one. For each $T \le 3$ the check computes exactly the ranks of $Q$ from ghost number $-1$ to $0$ and from $0$ to $1$, so the dimension of the cohomology at ghost number zero, and the signature of the Krein form on $\ker Q$; the results are $\dim H = 1, 2, 3, 4$ and signature $(\dim H, 0, \dim\operatorname{im}Q)$ for both momenta. Since $Q$ is self-adjoint and nilpotent, $\operatorname{im}Q \subset \ker Q$ is orthogonal to $\ker Q$, so the form descends to the cohomology, where it is positive definite by the signature. The pattern $T + 1$ is the dimension of the degree-$T$ symmetric power of a two-dimensional space; that it continues for every $T$ is the Kugo-Ojima quartet mechanism (Fact \claimref{synthesis.imported-brst-kugo-ojima-strocchi-and-gribov}(ii)), not computed here.

\emph{(iv).} At $T = 1$ and ghost number zero the space is spanned by the linear monomials $\epsilon\cdot z$; $Q(\epsilon\cdot z) = d^{*}L(\epsilon\cdot z)$ vanishes exactly when $k\cdot\epsilon = 0$, and the image of $Q$ from ghost number $-1$ is $L^{\dagger}$ applied to the ghost state, the monomial along $k$. So the one-particle cohomology is $k^{\perp}/\mathrm{span}\{k\}$, two-dimensional. For $\epsilon \in k^{\perp}$, $k\epsilon = k\wedge\epsilon$ is a bivector and $k(k\epsilon) = Q(k)\epsilon = 0$; $k\wedge\epsilon = 0$ exactly when $\epsilon$ is a multiple of $k$. So $\epsilon \mapsto k\epsilon$ induces an injective map from $k^{\perp}/\mathrm{span}\{k\}$ into $\{f : kf = 0\}$, which is two-dimensional (Proposition \claimref{field-theory.fock-space-maxwell-modes}(2), recomputed in the check as the kernel of $f \mapsto kf$ on the six bivectors), hence onto. The check confirms rank $2$ of the image and the dimension $2$ for three null momenta.

\emph{(v).} Write $T_a$ for the twelve anti-Hermitian generators of $\mathfrak{su}(2) \oplus \mathfrak{su}(3) \oplus \mathfrak{u}(1)$ on $\C^{2} \otimes \C^{3}$, the algebra of Theorem \claimref{standard-model.gauge-algebra-minimal-massive-composite} (a basis with Gaussian-rational entries; the structure constants $[T_a, T_b] = f_{ab}{}^{c}T_c$ are then rational and are computed exactly). Take $c^a = v_a$, $\bar c_a = w_a$ on twelve tower planes, so $\{c^a, \bar c_b\} = \delta^a_b$ and the $c$'s anticommute. With $Q = c^aT_a + G$, $G = -\tfrac12 f_{ab}{}^{c}c^ac^b\bar c_c$,
\[
Q^2 = \tfrac12 c^ac^b\,[T_a, T_b] + T_d\,\{c^d, G\} + G^2 ,
\]
since the matter factors commute with the ghosts. Moving $c^d$ through $c^ac^b$ costs no sign, so $\{c^d, c^ac^b\bar c_c\} = c^ac^b\{c^d, \bar c_c\} = \delta^d_c\,c^ac^b$ and $\{c^d, G\} = -\tfrac12 f_{ab}{}^{d}c^ac^b$. The first two terms therefore cancel exactly when $[T_a, T_b] = f_{ab}{}^{c}T_c$, and $G^2 = 0$ exactly when the $f$ satisfy the Jacobi identity (the standard computation of the BRST charge of a constraint algebra, valid for any Lie algebra and finite-dimensional representation; Fact \claimref{synthesis.imported-brst-kugo-ojima-strocchi-and-gribov}(i)). The check does not rely on this expansion: it computes $Q^2$ as a matrix of elements of the twelve-plane floor, finding it zero; dropping $G$ leaves $\tfrac12 c^ac^bf_{ab}{}^{c}T_c \neq 0$, and changing one structure constant of $\mathfrak{su}(3)$ by one makes $Q^2 \neq 0$.

\emph{Reading.} By (i) the store's representations carry no self-adjoint BRST charge as they stand, and by (ii) to (iv) the Gupta-Bleuler form of the polarisations supplies the missing indefinite form in a boost-invariant way, with physical states that are positive and coincide at one particle with the stored photon modes; under the premise that gauge fields carry this indefinite form, Premise \claimref{synthesis.krein-form-premise-indefinite-inner-product-for-gauge-fields}, the covariant route works at the free level. An indefinite form already appears at the one-particle level of the attachment, the twistor form of signature $(2,2)$ whose positive polarisations alone carry states (Theorem \claimref{synthesis.positive-structure-is-twistor-form-at-imaginary-time}); that is a parallel, not the form used here, and nothing is computed from it. The ghosts in (ii) and (v) are tower modes, as the Grassmann variables of Proposition \claimref{field-theory.path-integral-inside-tower} are rungs, and they are spin-zero CAR modes; Proposition \claimref{field-theory.twistor-attachment} makes a CAR lift local only for odd $2s$ when the weights are positive, and the spacelike bracket of a spin-zero CAR field vanishes only at antiparticle weight $-1$, that is in the indefinite form (\claimref{synthesis.locality-of-the-attached-free-fields-fixes-their-statistics}); they appear only in the zero-norm quartets of (iii). The adjoint representation of the derived algebra is real, so the gauge fields need no complex unit beyond the field's $I_4$ (argued). Not computed: the field-level charge of the interacting theory, the alternative physical (Coulomb) gauge for non-abelian fields with its Gribov ambiguity (Fact \claimref{synthesis.imported-brst-kugo-ojima-strocchi-and-gribov}(iv)), and confinement; the sign of the one-loop beta coefficients is computed in a result (\claimref{synthesis.one-loop-running-coefficients-from-spacetime-algebra-spin-traces}).
